\documentclass[10pt,a4paper]{article}
\usepackage[utf8]{inputenc}
\usepackage[francais]{babel}
\usepackage[T1]{fontenc}
\usepackage{amsmath}
\usepackage{amsfonts}
\usepackage{amssymb}
\usepackage{graphicx}
\usepackage[left=2cm, top=3cm, bottom=2cm, right=2cm]{geometry}
\usepackage{multirow}
\usepackage{lmodern}
\usepackage{tabularx}
\usepackage{colortbl}
\usepackage{pdflscape} % Package pour le mode paysage

\definecolor{lavender}{RGB}{230,230,250}

% Définition du compteur rownum
\newcounter{rownum}

\begin{document}

\begin{landscape} % Début du mode paysage
\begin{center}
\begin{tabularx}{\linewidth}{|c|c|X|m{2.5cm}|r|r|c|m{2cm}|}
    \hline
    \rowcolor{lavender} \textbf{N.} & \textbf{Numéro de Contrat} & \textbf{Nom du Client} & \textbf{Date Sign} & \textbf{Montant (FCFA)} & \textbf{Durée (mois)} & \textbf{Statut} & \textbf{État} \\ \hline
    \BLOCK{for contrat in contrats}
        \addtocounter{rownum}{1} % Incrémentation du compteur
        \therownum & CTR\VAR{contrat.id} & \VAR{contrat['client']['nom']} & \VAR{contrat['date_contrat']} & \VAR{contrat['montant_contrat']} & \VAR{contrat['duree']} & \VAR{contrat['status']} & \VAR{contrat['etat']}\\ \hline
    \BLOCK{endfor}
\end{tabularx}
\end{center}
\end{landscape} % Fin du mode paysage

\end{document}
